\documentclass[11pt,letterpaper]{article}
\usepackage[T1]{fontenc}
\usepackage{lmodern,microtype,amsmath,amssymb,amsthm,booktabs}
\usepackage[margin=1.04in]{geometry}
\usepackage{xcolor,hyperref,fancyhdr}
\definecolor{ink}{HTML}{134E58}
\hypersetup{colorlinks=true,linkcolor=ink,urlcolor=ink,citecolor=ink,pdftitle={Sharp localization at the zeros of completed Dirichlet L-functions}}
\pagestyle{fancy}\fancyhf{}
\fancyhead[L]{\small\textsc{Dirichlet zero filters}}
\fancyhead[R]{\small Sharp localization}\fancyfoot[C]{\thepage}
\setlength{\headheight}{14pt}\setlength{\emergencystretch}{2em}
\newtheorem{theorem}{Theorem}[section]
\newtheorem{lemma}[theorem]{Lemma}
\newtheorem{proposition}[theorem]{Proposition}
\newcommand{\R}{\mathbb R}\newcommand{\C}{\mathbb C}\newcommand{\E}{\mathbb E}
\newcommand{\dd}{\,d}\newcommand{\ord}{\operatorname{ord}}
\newcommand{\calF}{\mathcal F}\newcommand{\calI}{\mathcal I}
\newcommand{\bcrit}{b_{\mathrm c}}
\title{Sharp localization at the zeros of completed Dirichlet \texorpdfstring{$L$}{L}-functions}
\author{Research note}\date{27 November 2026}
\begin{document}
\maketitle
\begin{abstract}
A positive averaging measure can annihilate every noncentral zero of a completed Dirichlet product. How rapidly can its tails decay? For a conjugation-closed multiset of $d$ primitive characters with geometric mean conductor $Q_*$, we improve the upper bound on the coefficient in the critical double-exponential moment from $2\pi d/Q_*$ to $\pi d/Q_*$. At this coefficient the Laplace transform of every admissible law is the central-zero-removed completion times a polynomial of explicitly bounded degree. The bound is the exact supremum for every family whose completion has no real noncentral zero. The lower bound uses finite signed sums of arbitrarily small translates to produce a strictly positive density, even when the original kernel changes sign. For $d$ copies of the even quadratic character modulo five, the endpoint is unattainable for $d\le2$, has a unique law for $3\le d\le10$, and has infinitely many laws for $d\ge11$. The statements for the full complex zero divisor are unconditional. Their interpretation using only real Fourier frequencies uses the critical-line property of the chosen factors.
\end{abstract}

\section{The localization problem}
An average over translations can suppress an oscillation using only nonnegative weights. If $X$ has law $\mu$, then $\E e^{i\gamma(t+X)}=e^{i\gamma t}\E e^{i\gamma X}$. A zero of the Fourier transform at $\gamma$ cancels that oscillation, and a zero of order $k$ also cancels its first $k-1$ polynomial modulations. A concentrated law uses translations close to the observation point. How much concentration remains possible after imposing the entire infinite list of zeros?

For the noncentral zero ordinates of a family of $d$ primitive Dirichlet characters, the ProveIt article establishes, under the critical-line property, that the largest possible exponent $\alpha$ for a moment $\E\exp(b e^{\alpha|X|})<\infty$ with some $b>0$ is $2/d$. It asks for the sharp coefficient $b$ at that exponent. Its Jensen estimate gives $b\le2\pi d/Q_*$, with $Q_*$ the geometric mean of the conductors. Its construction provides a positive law for some $b>0$ by smoothing the inverse transform of the completed product. The modulus-one case uses Riemann's positive theta kernel. The repository presents the question, the upper bound and the smoothing argument in its section on positive filters \cite{proveit,riemann}.

We determine the exact supremum for every family without a real noncentral zero. It is one half of the earlier upper bound. In particular, for the zeta zeros, the range of attainable coefficients in $\E\exp(b e^{2|X|})<\infty$ is exactly $0<b<\pi$. The normalized Riemann kernel works simultaneously for every $b<\pi$, and no law works at $b=\pi$. For repeated copies of the even quadratic character modulo five, a law first realizes the supremum at three copies, then loses uniqueness at eleven copies. Both transitions come from a strict inequality on an integer degree.

\subsection{Completions and the prescribed divisor}
Let $\calF$ be a nonempty finite multiset of primitive Dirichlet characters closed under complex conjugation. Repetitions count in $d=|\calF|$. Write $a_j\in\{0,1\}$ for the parity of the $j$th character and $q_j$ for its conductor. For a nontrivial primitive character use
\[
 \Lambda_j(s)=(q_j/\pi)^{(s+a_j)/2}\Gamma((s+a_j)/2)L(s,\chi_j).
\]
For the modulus-one character put $q_j=1$, $a_j=0$ and use
$\xi(s)=\tfrac12s(s-1)\pi^{-s/2}\Gamma(s/2)\zeta(s)$.
Write $n_0$ for the number of copies of this last factor and set
\begin{gather}
 A=d/2,\qquad Q_*=(q_1\cdots q_d)^{1/d},\qquad \bcrit=\pi d/Q_*,
 \label{eq:constants}\\
 F(w)=\prod_{j=1}^d\Lambda_j(\tfrac12+w),\qquad m=\ord_0 F,\qquad
 H(w)=\frac{F(w)/w^m}{\lim_{v\to0}F(v)/v^m},\label{eq:H}\\
 \kappa=2n_0-m+\sum_{j=1}^d\left(\frac{a_j}{2}-\frac14\right).
 \label{eq:kappa}
\end{gather}
Here $\Lambda_j$ also denotes $\xi$ where appropriate. The functional equation and conjugate pairing make $F$ real on the real axis and give $F(-w)=\eta F(w)$ with $\eta=(-1)^m$, by the first nonzero Taylor coefficient. Thus $H$ is an even entire function of order at most one with $H(0)=1$. Removing the central zero is necessary since a probability transform has value one at zero. We use the classical analytic properties of primitive completions \cite{dirichlet,davenport}.

For a probability law write
\[
 M(w)=\int_{\R}e^{wt}\dd\mu(t),\qquad
 \calI_b(\mu)=\int_{\R}\exp\!\left(b e^{|t|/A}\right)\dd\mu(t).
\]
Call the law \emph{admissible} if its defining Laplace integral is entire and $M$ contains every zero of $H$ with at least its multiplicity. This prescribes the full complex divisor. If the nontrivial zeros of the chosen factors lie on $\Re s=1/2$, this asks precisely for the required zero of $M(i\gamma)$ at every nonzero real ordinate, including sums of multiplicities when factors share a zero. The full-divisor formulation and the proofs below also apply without that input. Competitors may be asymmetric.

\begin{theorem}[Critical coefficient and endpoint rigidity]\label{thm:main}
If $\mu$ is admissible and $\calI_b(\mu)<\infty$ for $b>0$, then $b\le\bcrit$. If $b=\bcrit$, there is a real polynomial $P$ such that
\begin{equation}
 M=PH,\qquad P(0)=1,\qquad \deg P+\kappa<-\tfrac12 . \label{eq:degree}
\end{equation}
Its degree is even, allowing degree zero. Conversely, if $PH$ is the entire Laplace transform of a probability law, $P$ is a polynomial and $P(0)=1$, the moment at $\bcrit$ is finite exactly under the strict degree inequality in \eqref{eq:degree}.
\end{theorem}
Existence of a positive inverse transform is part of the converse hypothesis. For an arbitrary family it is a genuine condition. For example, a real zero of $H$ makes admissibility impossible because $M(u)>0$ for real $u$. For the explicit positive kernels in Section~\ref{sec:examples} one can check the condition directly.

\section{Division by the completion}
Jensen's formula controls the circular mean of $\log|M|$. The quotient $M/H$ retains directional information. At the critical coefficient Stirling's formula yields a quadratic bound in the angle away from either real ray, on the scale of the exponential growth rate. A maximum principle transfers these estimates through the intervening sectors. The resulting quotient has zero exponential type. Its polynomial bound on the real axis then determines its algebraic form.

\begin{lemma}[Directional comparison]\label{lem:stirling}
If $\calI_b(\mu)<\infty$, $M$ is entire of order at most one. For an admissible law $G=M/H$ is also entire of order at most one. On a fixed right-hand ray $w=r e^{i\theta}$, $|\theta|<\pi/2$, let $c=\cos\theta$. Then
\begin{equation}
 \limsup_{r\to\infty}\frac{\log|G(re^{i\theta})|}{r}
 \le A\left[c\log\!\left(\frac{\bcrit}{b}c\right)+\theta\sin\theta\right].
 \label{eq:direction}
\end{equation}
The same bound holds on the opposite ray.
\end{lemma}
\begin{proof}
Set $v=|\Re w|$. Pull $\sup_{t\ge0}\exp(vt-b e^{t/A})$ out of the moment integral. For $Av\ge b$ this gives
\begin{equation}
 |M(w)|\le\calI_b(\mu)\exp\left[Av\left(\log\frac{Av}{b}-1\right)\right].
 \label{eq:moment}
\end{equation}
On compact sets the integrand and its derivatives admit integrable majorants, so $M$ is entire. Its growth is $\exp(O(r\log(2+r)))$. Admissibility removes all potential poles of $G$.

Write $m^+(r,f)$ for the circular mean of $\log^+|f|$. Jensen's formula and $H(0)=1$ give $m^+(r,1/H)\le m^+(r,H)$ at radii avoiding zeros. Hence $m^+(r,G)\le m^+(r,M)+m^+(r,H)$. Since $H$ has order at most one, the last sum is $O_\epsilon(r^{1+\epsilon})$ for every $\epsilon>0$. The Poisson majorant of the subharmonic function $\log^+|G|$ on a disk with radius between $2r$ and $3r$ bounds its maximum on $|w|\le r$ by three times this mean. This proves the order assertion.

On a closed sector $|\arg w|\le\pi/2-\delta$ the Dirichlet series gives $L(1/2+w,\chi_j)=1+O(2^{-\Re w})$ as $|w|\to\infty$. Stirling's expansion for $\Gamma(w/2+(1/2+a_j)/2)$ gives, with the principal logarithm and a nonzero real $C$,
\begin{equation}
 H(w)=C w^\kappa\exp\left[Aw\left(\log w-1+\log\frac{A}{\bcrit}\right)\right]
       (1+O_\delta(|w|^{-1})).
 \label{eq:stirling}
\end{equation}
Indeed $A/\bcrit=Q_* /(2\pi)$, each gamma factor contributes the power $a_j/2-1/4$, each $s(s-1)$ adds two powers, and division by $w^m$ removes $m$ powers. We use the usual sectorial Stirling expansion \cite{gamma}.

For $w=r e^{i\theta}$ the real part of $w\log w$ is $r(c\log r-\theta\sin\theta)$. Subtracting the logarithm of the modulus in \eqref{eq:stirling} from \eqref{eq:moment} cancels the terms of order $r\log r$ and gives \eqref{eq:direction}. On the opposite ray use evenness of $H$ and the same bound on $|M|$.
\end{proof}

\begin{lemma}[Transfer across a sector]\label{lem:sector}
Suppose $G$ is entire of order at most one. In a sector of opening $2\phi<\pi$ with middle direction $\psi$, suppose $|G(w)|\le C e^{k|w|}$ on both bounding rays, with real $k$. Then throughout the sector
\begin{equation}
 |G(w)|\le C'\exp\left(\frac{k}{\cos\phi}\Re(e^{-i\psi}w)\right).
 \label{eq:sector}
\end{equation}
\end{lemma}
\begin{proof}
Multiply by $\exp[-k e^{-i\psi}w/\cos\phi]$. The result is bounded on the rays and has order at most one. On truncated sectors choose $1<p<\pi/(2\phi)$ and multiply further by $\exp[-\epsilon(e^{-i\psi}w)^p]$, using the branch real on the middle ray. The real part of $(e^{-i\psi}w)^p$ is at least $|w|^p\cos(p\phi)>0$ away from the origin. Because the order is strictly below $p$, the circular boundary at infinity contributes nothing to the maximum principle. Let $\epsilon\downarrow0$. A fixed disk about the origin is absorbed in the constant.
\end{proof}

\begin{proof}[Critical bound and polynomial rigidity in Theorem~\ref{thm:main}]
The rays $\pm\pi/4$ and $\pi\pm\pi/4$, together with Lemmas~\ref{lem:stirling} and~\ref{lem:sector}, first give finite exponential type for $G$ at every $b$ under consideration. To refine this estimate use the four rays $\pm\theta,\pi\pm\theta$, with $0<\theta<\pi/4$. Their sectors have openings $2\theta$ and $\pi-2\theta$.

If $b>\bcrit$, the right side of \eqref{eq:direction} is negative at $\theta=0$ and stays negative for small $\theta$. Choose $k<0$ above its value for all four rays. Formula \eqref{eq:sector} bounds $G$ on the plane. Liouville's theorem makes $G$ constant, and its decay on a boundary ray forces it to vanish, contradicting $G(0)=1$.

At $b=\bcrit$ the right side of \eqref{eq:direction} is
$A(\cos\theta\log\cos\theta+\theta\sin\theta)=A\theta^2/2+O(\theta^4)$.
Choose $k=C_0\theta^2$ on all four rays for fixed $C_0>A$ and small $\theta$, allowing the multiplicative constant to depend on $\theta$. In the large sectors the denominator $\cos\phi$ in \eqref{eq:sector} equals $\sin\theta$, and in the small ones it equals $\cos\theta$. Thus for every $\epsilon>0$ one has $|G(w)|\le C_\epsilon e^{\epsilon|w|}$ on the whole plane.

Equations \eqref{eq:moment} and \eqref{eq:stirling} on the real axis give $G(u)=O(|u|^{-\kappa})$ as $u\to\pm\infty$. Choose a nonnegative integer $N\ge-\kappa$. In the upper half-plane $g(w)=G(w)/(w+i)^N$ has zero exponential type and is bounded by $C_1\ge1$ on the real axis. In the upper half-disk of radius $R$ the harmonic measure of the semicircle at fixed $w$ is
\[
 \frac{2}{\pi}\arg\frac{R+w}{R-w}=O_w(R^{-1}).
\]
Since $\sup_{|z|=R,\ \Im z\ge0}\log^+|g(z)|=o(R)$, the two-constants estimate gives $\log|g(w)|\le\log C_1+o(1)$. In the lower half-plane divide instead by $(w-i)^N$. Thus $|G(w)|=O((1+|w|)^N)$ on the plane, and Cauchy's estimates make $G$ a polynomial $P$. The functions $M,H$ are real on the real axis and $H(0)=1$, so $P$ has real coefficients and $P(0)=1$. Evenness of $H$ and positivity of $M$ imply that $P$ has the same sign at both real infinities. Its degree is even. The next calculation yields the exact degree bound.
\end{proof}

\section{The last half power}
For a positive law with all real exponential moments, Tonelli's theorem gives
\begin{equation}
 S_b:=\sum_{n=0}^{\infty}\frac{b^n}{n!}\big(M(n/A)+M(-n/A)\big),
 \qquad \calI_b(\mu)\le S_b\le2\calI_b(\mu).
 \label{eq:series}
\end{equation}
Indeed, after summation the maximum of the two exponential integrands is $\exp(b e^{|t|/A})$. Positivity of every summand permits an exact criterion even at the threshold of summability.

Suppose $M=PH$ with a polynomial of degree $k$. Since $M>0$ on both real rays and $H$ is even, $k$ is even and the leading coefficients add in $M(u)+M(-u)$. At $u=n/A$, formula \eqref{eq:stirling} and Stirling's formula for $n!$ give
\begin{equation}
 \frac{b^n}{n!}\big(M(n/A)+M(-n/A)\big)
   \sim C_P\left(\frac{b}{\bcrit}\right)^n n^{\kappa+k-1/2},
   \qquad C_P>0 .
 \label{eq:p-series}
\end{equation}
At $b=\bcrit$ this series converges exactly when $\kappa+k<-1/2$. In the equality case the partial sums diverge logarithmically. This finishes both implications of Theorem~\ref{thm:main}. It also shows that whenever $H$ is itself a probability transform, a single law realizes every $b<\bcrit$.

\section{Sharpness for general Dirichlet products}\label{sec:general}

A signed combination of small translations can repair a sign-changing kernel while preserving its zeros. The coefficients may have either sign. The resulting density is strictly positive. The arithmetic input is the gap between the first and second terms of each ordinary Dirichlet series.

\begin{theorem}[Sharp supremum for all admissible families]\label{thm:general}
Suppose the central-zero-removed product $H$ has no real zero. For every $0<b<\bcrit$ there is an admissible strictly positive continuous probability density with $\calI_b<\infty$. Hence $\bcrit$ is the exact supremum for this family. If $H$ has a real zero, no admissible probability law exists. In particular the upper bound is sharp as a supremum for every family satisfying the critical-line property.
\end{theorem}

We give the proof in three steps. Throughout this section assume $H$ has no real zero. Let $h$ be its real even inverse bilateral Laplace kernel. The repository constructs this continuous kernel by convolving the theta kernels and integrating $m$ times from the right to remove the central zero. It also proves that there exists a nonnegative compactly supported integrable $k_0$ such that $h*k_0>0$ everywhere. Indeed $H(u)>0$ on $\R$, $h$ is even, positive on both tails, and has every exponential moment in total variation. The signed-to-positive smoothing lemma in \cite{proveit} applies. We will use this single broad smoothing only at the end of a separation argument.

\subsection*{The first Dirichlet coefficient}
For $c_j=1/2+a_j$, let
\[
 v_j(t)=2e^{c_jt}\exp(-\pi e^{2t}/q_j),\quad
 v=v_1*\cdots*v_d,\quad \beta_0=d^{-1}\sum_j c_j .
\]
Let $Rg(t)=\int_t^\infty g(u)\dd u$, put $v_m=R^m v$ with $R^0$ the identity,
and let $D$ denote differentiation in $t$. There is a positive constant $c_*$ for which
\begin{equation}
 h(t)=c_*\sum_{n\ge1}\frac{\alpha_n}{\sqrt n}\,
             B(D)v_m(t+\log n),\quad
 B(D)=\left(\frac{D^2-1/4}{2}\right)^{n_0}.                 \label{eq:Dir-exp}
\end{equation}
Here $\alpha_1=1$ and $\sum_n\alpha_n n^{-s}=\prod_jL(s,\chi_j)$ for $\Re s>1$. In particular $|\alpha_n|$ grows at most polynomially. To verify the formula, the bilateral transform of $v_j$ is the conductor and gamma part of the $j$th completion. Integration from the right divides its transform by $w$ for real $w>0$, and $B(D)$ supplies the zeta factors $s(s-1)/2$. Translation by $\log n$ supplies $n^{-w}$. These calculations and the sum are absolutely integrable on a sufficiently far right vertical line. For example, each $v_j$ and each of its fixed derivatives has an integrable absolute exponential tilt there, as does $R^m v$, by Tonelli's theorem applied to $|v|$. Derivatives commute with $R$ here by decay at $+\infty$. Fourier uniqueness on that line gives \eqref{eq:Dir-exp}. The normalization $c_*$ is positive since $F(w)/w^m$ is positive for large positive $w$ and has no real zero after central removal.

There is a positive measure $\rho$ supported on $[\bcrit,\infty)$ with infimum of support exactly $\bcrit$ and
\begin{equation}
 v_m(t)=e^{\beta_0 t}\int_{\bcrit}^\infty e^{-u e^{t/A}}\dd\rho(u).
 \label{eq:leading}
\end{equation}
In the integral for $v$, write $t_j=t/d+y_j$ with $\sum_jy_j=0$. The exponent is $-e^{t/A}\sum_j(\pi/q_j)e^{2y_j}$. The positive weight on this hyperplane is a constant times $\exp(\sum_j c_j y_j)\dd y_1\cdots\dd y_{d-1}$. The sum in the exponent attains its minimum $\bcrit$ when all $(\pi/q_j)e^{2y_j}$ coincide; this is the arithmetic--geometric mean inequality. The pushforward of the weight has positive mass arbitrarily close to the minimum. For $d=1$ the pushforward is a point mass. For $m>0$ use
$R^mv(t)=\int_0^\infty z^{m-1}v(t+z)\dd z/(m-1)!$.
This mixes the same positive measure under the dilations $u\mapsto u e^{z/A}$ with positive weights. Its infimum of support stays $\bcrit$. The integral in \eqref{eq:leading} is finite for every $e^{t/A}>0$.

Here are two consequences that control the convergence boundaries below. On the right tail,
\begin{equation}
 h(t)>0,\qquad
 \lim_{t\to+\infty} e^{-t/A}\log h(t)=-\bcrit .             \label{eq:exacttail}
\end{equation}
Indeed put $x=e^{t/A}$. Applying $B(D)$ under \eqref{eq:leading} replaces $e^{-xu}$ by $e^{-xu}$ times a polynomial in $xu$ of degree $2n_0$ with positive leading coefficient, since $D$ acts as $\beta_0+(x/A)\partial_x$ after the prefactor is removed. For $x$ large this polynomial is positive uniformly on $u\ge\bcrit$ and bounded below by a positive constant. Mass of $\rho$ on $[\bcrit,\bcrit+\epsilon]$ gives the lower logarithmic rate, and exponential damping of every polynomial gives the upper one. In \eqref{eq:Dir-exp} all $n\ge2$ have the strictly faster rate coming from $n^{1/A}x$. More precisely, for complex $s$ in a compact subset of $\Re(e^{-s/A})>0$ and large real $t$, the absolute sum over $n\ge J$ in $h(t-s)$ is bounded by
\begin{equation}
 C\exp\big[-(\bcrit-\epsilon)J^{1/A}\Re(e^{-s/A}) e^{t/A}\big],
 \qquad J=1,2,                                               \label{eq:tailanalytic}
\end{equation}
for every fixed $0<\epsilon<\bcrit$. To see uniformity, each integrand has a fixed polynomial factor in $n^{1/A}xu$ and a power of $n$ from $\alpha_n n^{\beta_0-1/2}$. A fraction of the exponential absorbs these factors and $e^{\beta_0 t}$. A further arbitrarily small fraction makes the $n$-sum bounded and supplies an integrable factor $e^{-u}$ against $\rho$. The compact lower bound for $\Re(e^{-s/A})$ makes these estimates uniform. This proves \eqref{eq:tailanalytic} by a convergent holomorphic sum, and proves \eqref{eq:exacttail} because $(\bcrit-\epsilon)2^{1/A}>\bcrit$ for small enough $\epsilon$. Evenness supplies the left tail.

\subsection*{Annihilating small translations}
Fix $\delta>0$. On the two tails choose $W(t)=h(|t|-\delta)$, extending it to a strictly positive continuous function on $\R$. By \eqref{eq:exacttail} the functions
\[
       f_s(t)=h(t-s)/W(t),\qquad |s|\le\delta ,
\]
are continuous on the compact interval $[-\infty,+\infty]$ after setting their endpoint values by limits. They vanish at both endpoints for $|s|<\delta$. The value at $+\infty$ of $f_\delta$ is one and its value at $-\infty$ is zero. For $f_{-\delta}$ the two values are interchanged.

We claim their real linear span contains a strictly positive function on this compact interval. Otherwise separation of the span from the open cone of strictly positive continuous functions, followed by the Riesz representation theorem, supplies a nonzero positive finite measure $\sigma$ that annihilates the span. On $\R$ set $\dd\lambda=\dd\sigma/W$. For $|s|<\delta$ this implies
\begin{equation}
                 \int_\R h(t-s)\dd\lambda(t)=0 .             \label{eq:orth}
\end{equation}
The measure $\lambda$ is positive and locally finite, but its tails can grow very rapidly. Define
\[
 r_\pm=\inf\{r>0:\int_{\{\pm t\ge0\}}\!
          \exp(-r e^{|t|/A})\dd\lambda(t)<\infty\},\qquad
 a_\pm=A\log(\bcrit/r_\pm),
\]
with $a_\pm=\infty$ for $r_\pm=0$. Formula \eqref{eq:exacttail} applied to $W$ shows $r_\pm\le\bcrit e^{-\delta/A}$, so $a_\pm\ge\delta$. By \eqref{eq:tailanalytic} and evenness, the left side of \eqref{eq:orth} is real analytic for $-a_-<s<a_+$, with a holomorphic neighborhood of every real point there. It vanishes throughout this interval by analytic continuation.

Suppose $a_+<\infty$, so $r_+>0$. Near $s=a_+$ the contribution of $t<0$ continues holomorphically since the opposite convergence boundary is at most $-\delta$. The contribution of the $n\ge2$ terms from $t\ge0$ also continues holomorphically by \eqref{eq:tailanalytic}, because
$2^{1/A}\bcrit e^{-a_+/A}>r_+$.
Bounded ranges of $t$ cause no issue by \eqref{eq:leading} and local uniform convergence. Thus the $n=1$ contribution from $t\ge0$ has analytic continuation across $a_+$. It equals $c_*B(-\partial_s)V(s)$ on the original side, where
\[
 V(s)=\int_{t\ge0}v_m(t-s)\dd\lambda(t).
\]
An ordinary differential equation with constant coefficients and nonzero leading coefficient has analytic solutions across an ordinary point whenever its right side continues analytically. Continuing from a nearby interior point shows that $V$ itself continues holomorphically through $a_+$. For $n_0=0$ the operator is the identity.

On the other hand, \eqref{eq:leading} gives a positive measure $\nu$ with
\[
 V(s)=e^{-\beta_0s}\int e^{-zv}\dd\nu(v),\qquad z=e^{-s/A}.
\]
Here $\nu$ is the pushforward under $v=u e^{t/A}$ of
$e^{\beta_0 t}\mathbf1_{t\ge0}\dd\lambda(t)\dd\rho(u)$. Its real Laplace abscissa is exactly $z_0=r_+/\bcrit$. For $z>z_0$, use $u\ge\bcrit$ and the small slack in $z\bcrit>r_+$ to absorb the prefactor and integrate over $\rho$. For $0<z<z_0$, choose $\epsilon>0$ with $z(\bcrit+\epsilon)<r_+$. The strictly positive $\rho$-mass below $\bcrit+\epsilon$ forces divergence, since $e^{\beta_0t}\ge1$. A positive Laplace integral cannot continue holomorphically across a finite real abscissa. Indeed, if it did, take $z_1>z_0$ sufficiently close that its Taylor disk extends to a real $z_2<z_0$. Expanding at $z_1$ towards $z_2$ gives nonnegative Taylor terms, and monotone convergence makes their sum equal to the divergent integral at $z_2$. This contradicts the continuation just proved, because $s\mapsto e^{-s/A}$ is locally invertible at $a_+$. Hence $r_+=0$. Applying the same argument after reflection gives $r_-=0$.

It follows that \eqref{eq:orth} holds for every real $s$. Its integrals converge absolutely and uniformly for $s$ in bounded sets by \eqref{eq:tailanalytic}. Integrate the identity against the compact kernel $k_0(s)\dd s$ from the start of the proof. Fubini's theorem gives $\int(h*k_0)(t)\dd\lambda(t)=0$. Strict positivity of $h*k_0$ forces $\lambda=0$. The remaining masses of $\sigma$ at $+\infty$ and $-\infty$ vanish by pairing respectively with $f_\delta$ and $f_{-\delta}$. This contradicts $\sigma\ne0$, proving the claim about the span.

\subsection*{Construction of a positive law}
Some finite real combination $g(t)=\sum_\ell c_\ell h(t-s_\ell)$, $|s_\ell|\le\delta$, therefore has $g/W$ strictly positive and bounded on the compactification. Normalize $g$ to integral one. Its bilateral transform is $H(w)$ times a finite sum of real multiples of $e^{s_\ell w}$, so it retains all prescribed zeros with their multiplicities. The upper bound $g\le C W$ and \eqref{eq:exacttail} give $\calI_b<\infty$ for every $b<\bcrit e^{-\delta/A}$. For any specified $b<\bcrit$ choose $\delta<A\log(\bcrit/b)$. This proves Theorem~\ref{thm:general}. The separation argument proves existence and gives no uniform bound on the number or absolute sizes of the coefficients.


\section{Positive kernels and exact constants}\label{sec:examples}
Put $\beta(s)=L(s,\chi_{-4})$, and let $\chi_5(n)=(n/5)$ denote the primitive quadratic character modulo five. For either nontrivial character in this list the theta kernel in logarithmic coordinates is
\begin{equation}
 K_\chi(t)=2e^{(1/2+a)t}\sum_{n\ge1}\chi(n)n^a\exp(-\pi n^2e^{2t}/q).
 \label{eq:theta}
\end{equation}
The bilateral Laplace transform is $\Lambda_\chi(1/2+w)$, by termwise Mellin integration for $\Re w$ large and analytic continuation. Specifically $x=e^{2t}$ gives the conductor and gamma factors in \eqref{eq:H}. Poisson summation for the primitive theta series gives $K_\chi(-t)=K_\chi(t)$ in both cases. The Gauss sums are $2i$ and $\sqrt5$, respectively, giving root number one after division by $i^a\sqrt q$. Theta inversion also supplies double-exponential decay on the left from the series on the right.

For $\chi_{-4}$ the nonzero terms alternate when $t\ge0$, starting positive. Their magnitudes decrease since $x\exp(-\pi x^2 e^{2t}/4)$ decreases for $x\ge1$. For $\chi_5$, put $c=\pi e^{2t}/5\ge\pi/5$. Relative to the first summand, the absolute sum of all later terms is at most
\[
 \sum_{n=2}^{\infty}e^{-c(n^2-1)}\le\frac{e^{-3\pi/5}}{1-e^{-\pi}}<1,
\]
because $n^2-1\ge3+5(n-2)$ for $n\ge2$. Thus both kernels are strictly positive for all real $t$.

For $\zeta$ use the Riemann kernel, given for $t\ge0$ by
\begin{equation}
 K_0(t)=\sum_{n\ge1}\left(4\pi^2 n^4 e^{9t/2}-6\pi n^2 e^{5t/2}\right)
                                e^{-\pi n^2 e^{2t}} .
 \label{eq:riemann}
\end{equation}
It is even by theta inversion and positive since $2\pi n^2e^{2t}>3$ on $t\ge0$. Its bilateral Laplace transform is $\xi(1/2+w)$, obtained by differentiating the Jacobi theta transformation, or by integrating the displayed two powers and using the gamma recurrence \cite{riemann,proveit}. Each of our three kernels has positive integral and hence no central zero. Divide each by its integral to obtain the probability density $p_j$.

\begin{proposition}[Sharp supremum for positive products]\label{prop:product}
For any nonempty product of these three completions, the exact supremum of $b$ is $\bcrit$. The convolution $p_1*\cdots*p_d$ realizes every $b<\bcrit$ and realizes $b=\bcrit$ exactly when $\kappa<-1/2$.
\end{proposition}
\begin{proof}
The transform of the convolution is $H$. Formula \eqref{eq:p-series} at $k=0$ proves its moment assertions, and Theorem~\ref{thm:main} excludes every $b>\bcrit$.

A direct estimate also explains the geometric mean conductor. Take independent $X_j$ with these densities. Equations \eqref{eq:theta} and \eqref{eq:riemann} imply $\E\exp(c_j e^{2|X_j|})<\infty$ for each $c_j<\pi/q_j$. The triangle inequality and the arithmetic--geometric mean inequality give
\[
 \frac{d}{Q_*}\exp\!\left(\frac2d\left|\sum_j X_j\right|\right)
 \le d\left(\prod_j\frac{e^{2|X_j|}}{q_j}\right)^{1/d}
 \le \sum_j\frac{e^{2|X_j|}}{q_j}.
\]
Multiply by $bQ_*/d<\pi$ and take exponential moments. Independence factors the last expectation into a finite product.
\end{proof}
Consequently the single zeta factor has exact supremum $\pi$, and the single beta factor has exact supremum $\pi/4$. Both are unattained because their respective values of $\kappa$ are $7/4$ and $1/4$. For one copy of each, the supremum is $\pi$, also unattained.

\section{Two changes of behavior modulo five}
Let $p$ denote the normalized kernel of $\chi_5$, and write $H_5$ for its one-copy Laplace transform. For the product of $d$ copies, $\bcrit=\pi d/5$, $\kappa=-d/4$, and the prescribed multiplicity at each noncentral zero grows by a factor of $d$.

\begin{theorem}[Attainment and uniqueness for every multiplicity]\label{thm:five}
For $d$ copies of $L(s,\chi_5)$ the following alternatives hold among all admissible probability laws.
\begin{center}
\begin{tabular}{@{}ll@{}}\toprule
Copies & Laws with $\calI_{\pi d/5}(\mu)<\infty$ \\ \midrule
$1\le d\le2$ & None.\\
$3\le d\le10$ & Exactly one, with density $p^{*d}$.\\
$d\ge11$ & Infinitely many, including $p^{*d}$.\\ \bottomrule
\end{tabular}
\end{center}
Every endpoint transform is $P H_5^d$, with $\deg P<d/4-1/2$ and even degree. Every nonnegative even integer degree satisfying this condition occurs.
\end{theorem}
\begin{proof}
Degree zero is allowed exactly when $d>2$. The next possible degree, two, is allowed exactly when $d>10$. Proposition~\ref{prop:product} gives the degree-zero law at all $d>2$. For $3\le d\le10$, \eqref{eq:degree} forces $P=1$ by normalization. Fourier uniqueness then identifies the law.

For each integer $j\ge1$ take $p_\epsilon=p+\epsilon p^{(2j)}$ with $\epsilon>0$ small. We show it is a positive probability density. With $x=\pi e^{2t}/5$ on the right half-line, derivatives of the first term in \eqref{eq:theta} have the form $C e^{t/2}e^{-x}R_k(x)$, where $R_0=1$ and
\[
 R_{k+1}(x)=(\tfrac12-2x)R_k(x)+2xR'_k(x).
\]
Thus $R_k$ has degree $k$ and leading coefficient $(-2)^k$. For fixed $k$ the other differentiated summands contribute $O_k(e^{t/2}x^k e^{-4x})$ as $x\to\infty$, by summing the polynomial factors in $n$ against $e^{-n^2x}$. Consequently $p^{(2j)}>0$ for large positive $t$, hence also for large negative $t$ by evenness. On the remaining compact interval $p$ has a positive minimum and $p^{(2j)}$ is bounded. This proves positivity for all sufficiently small $\epsilon$. The tail estimates give $\int p^{(2j)}=0$ and justify integration by parts against $e^{wt}$. The Laplace transform of $p_\epsilon$ is $(1+\epsilon w^{2j})H_5(w)$.

Convolve $p_\epsilon$ with $p^{*(d-1)}$, using the point mass convention for the convolution power zero. By \eqref{eq:p-series} this law has the endpoint moment exactly when $2j<d/4-1/2$. Different $\epsilon$ give different transforms. Taking $j=1$ proves the third line of the table, and varying $j$ proves the degree assertion.
\end{proof}
The first new degree $2j$ appears at $d=8j+3$. With two copies the canonical series has terms asymptotic to a positive multiple of $n^{-1}$. At ten copies the same logarithmic divergence excludes degree two. The uniqueness proof covers competitors with arbitrary additional zeros, regardless of whether their kernels are smooth or even.

\section{Scope and further questions}
The gamma factors encode the conductor-dependent coefficient in \eqref{eq:stirling}. The factors $s(s-1)$, parity and the central multiplicity survive as the power $\kappa$. At the critical coefficient the directional estimate forces a polynomial quotient, and the moment series detects the remaining half power.

A classical structural analogy is Hardy's theorem on Gaussian Fourier decay, whose critical hypotheses also leave a finite-dimensional space of polynomial modifications \cite{hardy}. Here the hypotheses prescribe the Dirichlet zero divisor and an integral moment on logarithmic scale. The claims of extension are relative to the pinned source in \cite{proveit}. We have not located the critical coefficient together with the polynomial degree criterion or the modulo-five attainment transitions in the earlier work consulted. A priority determination would require a wider review.

Two questions remain open in this treatment.
\begin{enumerate}
\item For a sign-changing inverse transform, which polynomials of the endpoint degree bound make its inverse transform nonnegative? In the subcritical construction, can one bound the number of translations and their coefficients as the target coefficient tends to $\bcrit$?
\item For the modulo-five family beyond ten copies, which normalized polynomials of admissible degree give a nonnegative inverse transform? Theorem~\ref{thm:five} realizes every allowed degree. A complete description would identify the boundary in coefficient space at which the perturbed convolution first has a real zero.
\end{enumerate}
The result determines the sharp supremum in the article's full-spectrum endpoint question for every conjugation-closed family under its critical-line hypothesis. A separate longstanding conjecture about the location or multiplicity of individual Dirichlet zeros would require further arithmetic input.

\IfFileExists{weighted.tex}{\section{A sharp endpoint refinement for the Riemann zeros}\label{sec:weighted}
For a single zeta factor the moment at $b=\pi$ always diverges. A linear term in the exponent detects exactly where the canonical law is forced. For an admissible law put
\[
  J_\lambda(\mu)=\int_{\R}
          \exp\!\left(\pi e^{2|t|}-\lambda |t|\right)\dd\mu(t),
          \qquad \lambda\in\R .
\]
\begin{theorem}
There is no admissible law with $J_\lambda<\infty$ for $\lambda\le9/2$. For $9/2<\lambda\le17/2$ there is exactly one, the normalized Riemann theta kernel. For $\lambda>17/2$ there are infinitely many.
\end{theorem}
\begin{proof}
Here $A=1/2$, $\kappa=7/4$, and $H(w)=\xi(1/2+w)/\xi(1/2)$. Suppose $J_\lambda<\infty$. The majorant used for \eqref{eq:moment} now has $v+\lambda$ in place of $v$, since the maximized expression is $(v+\lambda)t-\pi e^{2t}$. For $v\to\infty$ this changes the logarithm of the bound by $A\lambda\log v+O_\lambda(1)$. Entire order and the indicator bound \eqref{eq:direction} at $b=\pi$ therefore persist. On either real ray the quotient $M/H$ is $O((1+|w|)^{A\lambda-\kappa})$ as $|w|\to\infty$. The zero-type and polynomial conclusions in the proof of Theorem~\ref{thm:main} apply without change. Let $k$ be the even degree of the resulting polynomial.

Expand $e^{\pi e^{2|t|}}$ by its positive series. For $n$ large enough that $u_n=2n-\lambda\ge0$, the moment $\int e^{u_n |t|}\dd\mu$ lies between $(M(u_n)+M(-u_n))/2$ and $M(u_n)+M(-u_n)$. The finitely many earlier terms are finite. Formula \eqref{eq:stirling} gives
\[
 \frac{H(2n-\lambda)}{H(2n)}
      \sim (n/\pi)^{-\lambda/2} .
\]
The summands of $J_\lambda$ are consequently bounded above and below by positive multiples of $n^{k+7/4-\lambda/2-1/2}$ for all large $n$. Hence convergence is equivalent to $\lambda>9/2+2k$ for every polynomial candidate that is a probability transform. Degree zero forces $M=H$ and the normalized positive kernel realizes it. Degree two first becomes possible above $17/2$.

For existence of degree two let $h$ be the normalized Riemann kernel. On the positive tail its defining series gives $h(t)\sim C e^{9t/2}e^{-\pi e^{2t}}$ and $h''(t)/h(t)\sim4\pi^2 e^{4t}$. These expansions follow by differentiating the summands: terms with index at least two are exponentially smaller than the first, even with the polynomial factors from two derivatives. The negative tail follows by evenness. On a compact interval $h$ has positive minimum and $h''$ is bounded. Thus $h+\epsilon h''$ is a probability density for every sufficiently small $\epsilon>0$, with transform $(1+\epsilon w^2)H(w)$ by integration by parts. These distinct laws satisfy $J_\lambda<\infty$ whenever $\lambda>17/2$.
\end{proof}
}{}
\begin{thebibliography}{9}
\bibitem{proveit} V. Reshetnikov, \emph{ProveIt}, Analysis/GeneralizedRiemannHypothesis, article section on positive zero filters. Commit \texttt{ff96fd44ab16}, file \texttt{fil.tex}, especially the results labelled \texttt{fil:thm:global}, \texttt{fil:lem:tails}, \texttt{fil:cor:endpoint}, and the open questions. \href{https://github.com/VladimirReshetnikov/ProveIt/blob/ff96fd44ab160e65fb8a9d861243828d07c156cf/Analysis/GeneralizedRiemannHypothesis/Article/sections/fil.tex}{Pinned source}.
\bibitem{riemann} B. Riemann, \emph{Ueber die Anzahl der Primzahlen unter einer gegebenen Gr\"osse}, Monatsberichte der Berliner Akademie (1859), 671--680. \href{https://www.claymath.org/collections/riemanns-1859-manuscript/}{Manuscript}.
\bibitem{dirichlet} NIST Digital Library of Mathematical Functions, \S25.15, \emph{Dirichlet $L$-functions}. \url{https://dlmf.nist.gov/25.15}
\bibitem{davenport} H. Davenport, \emph{Multiplicative Number Theory}, 3rd ed., revised by H. L. Montgomery, Graduate Texts in Mathematics 74, Springer, 2000.
\bibitem{gamma} NIST Digital Library of Mathematical Functions, \S5.11(i), \emph{Poincar\'e-type expansions}. \url{https://dlmf.nist.gov/5.11}
\bibitem{hardy} G. H. Hardy, \emph{A theorem concerning Fourier transforms}, J. London Math. Soc. s1-8 (1933), 227--231. \href{https://doi.org/10.1112/jlms/s1-8.3.227}{doi:10.1112/jlms/s1-8.3.227}.
\end{thebibliography}
\end{document}
